\section{Related Work}
\label{sec:related}

\paragraph{Speculative decoding.}
Draft-then-verify decoding goes back to blockwise parallel decoding~\cite{stern2018blockwise}.
Speculative decoding and speculative sampling~\cite{leviathan2023fast,chen2023accelerating} added an acceptance rule that preserves the target distribution, so under greedy decoding the output equals the target's own argmax sequence, in exact arithmetic.
Later work improves the drafter: extra decoding heads (Medusa~\cite{cai2024medusa}), feature-level drafters (EAGLE~\cite{li2024eagle,li2024eagle2,li2025eagle3}), a recurrent drafter (ReDrafter~\cite{cheng2024redrafter}), layer skipping~\cite{zhang2024draftverify}, Jacobi-style n-gram guessing~\cite{fu2024lookahead}, and drafters that emit a whole block in one pass (DFlash~\cite{chen2026dflash}, DFlare~\cite{zhang2026dflare}, Liquid AI's DSpark~\cite{liquid2026dspark}).
Others verify token trees instead of chains~\cite{miao2024specinfer,li2024eagle2,xiong2024dyspec}.
FR-Spec~\cite{zhao2025frspec} restricts the draft LM head to a frequency-ranked vocabulary subset, and the released EAGLE-3 checkpoints draft over 32k of 128k tokens; our 43.8k-token MTP head applies this idea to Qwen3.5's 248k vocabulary and is not new.
Adaptive draft length (SpecDec++~\cite{huang2025specdecpp}, PEARL~\cite{liu2025pearl}) is likewise established, as is overlapping drafting with verification: PEARL drafts during verification, and Speculative Speculative Decoding~\cite{kumar2026ssd} prepares speculations for predicted verification outcomes.
Our optimistic pipelining is the single-outcome case of that idea (predict full acceptance) on one device.
What we add there is an MLX-specific hazard: in-place KV-cache writes inside a discarded lazy graph still execute unless the KV views are pinned.
Our negative result for tree drafting reflects our measured verification cost curve; tree and block methods assume verification of many tokens is nearly free, which it is not on an M1 Pro.

\paragraph{Multi-token prediction.}
Gloeckle et al.~\cite{gloeckle2024mtp} train several output heads on a shared trunk and report up to $3\times$ faster inference by using them for self-speculation.
DeepSeek-V3~\cite{deepseek2024v3} trains sequential MTP modules and reuses them as a drafter, with 85--90\% second-token acceptance and $1.8\times$ tokens per second.
Qwen3-Next~\cite{qwen2025qwen3next} and Qwen3.5~\cite{qwen2026qwen35} ship MTP modules intended for inference; Qwen3.5-9B stacks eight blocks of three Gated DeltaNet layers and one gated-attention layer~\cite{qiu2025gatedattn}.
MTP self-speculation for this family already exists in llama.cpp~\cite{am17an2026llamacppmtp}, in an open mlx-lm pull request~\cite{airrunner2026mlxlmmtp}, and in MTPLX~\cite{altoukhi2026mtplx}.
We claim no novelty for MTP drafting.
Our contribution is making its verification cheap and exact on hardware where, by default, it is neither.

\paragraph{Linear-attention and hybrid models.}
Gated DeltaNet~\cite{yang2025gdn} combines the delta rule of DeltaNet~\cite{yang2024deltanet} with the gating of Mamba-2~\cite{dao2024mamba2,gu2023mamba}.
Its layers overwrite a fixed-size state, so rejected drafts cannot be dropped by truncating a cache.
The Mamba in the Llama~\cite{wang2024mambainllama} gave a speculative decoding algorithm for distilled hybrids and STree~\cite{wu2025stree} gave tree verification for SSMs.
Work from 2026 targets Gated DeltaNet hybrids in GPU serving engines: TreeWY~\cite{ghantasala2026treewy} and Bole~\cite{wang2026bole} compute tree updates in closed form and rebuild only the accepted state, SpecLA~\cite{wang2026specla} stores compact factors to recover it, Weaver~\cite{oda2026weaver} verifies trees without rollback, KVBuffer~\cite{zou2026kvbuffer} defers state updates by buffering keys and values, LumoTree~\cite{ma2026lumotree} (first released as ``GDN Tree-Scan'') replays branch-local scans inside vLLM, and ReplaySSM~\cite{dao2026replayssm} caches inputs instead of states.
TreeWY and ReplaySSM both describe snapshotting the recurrent state at every draft position as current practice~\cite{ghantasala2026treewy,dao2026replayssm}, and the mlx-lm pull request does the same~\cite{airrunner2026mlxlmmtp}.
Our rollback steps the recurrence per position inside the batched verify pass and keeps every (conv, SSM) state, so it is that standard technique, not a new one.
It fits our setting because one sequence, chains of at most four tokens, and a 9B model make the snapshots small, whereas the systems above attack the memory and tree-width limits of batched GPU serving.

\paragraph{On-device and Apple Silicon.}
ReDrafter~\cite{cheng2024redrafter} reported up to $2.3\times$ on Apple Silicon with an implementation in MLX~\cite{hannun2023mlx}.
Chordiya~\cite{chordiya2026lossless} found on an M3-class laptop that the quantized Metal backend of llama.cpp~\cite{gerganov2023llamacpp} verifies draft tokens serially, with verify cost growing nearly linearly in the number of drafts for small batches, and three of five configurations ran slower than plain decoding.
We see the same effect in MLX 0.32.4 on an M1 Pro.
MLX's small-batch kernel \texttt{qmv\_wide}~\cite{gross2026qmvwide}, adapted from llama.cpp's \texttt{mul\_mv\_ext}, is enabled for affine quantization only on GPU generation 15 (M3) and newer, so 4-bit verification on M1-class GPUs still runs one mat-vec per token.
llama.cpp has since added few-row simdgroup-matrix kernels for speculative verification on Apple7+ GPUs~\cite{narola2026fewrow} (build b11404), after the b10969 build we use as a baseline.
Our multi-token kernel uses the same idea, reading each weight once for $T$ tokens, for affine 4-bit weights on an M1 Pro; what sets it apart is the summation-order constraint described next.
MTPLX reports $1.6\times$ on a 16\,GB M4 Mac mini and the mlx-lm pull request $1.57\times$ for a 27B model on an M4 Pro; these figures use other chips and models and are not directly comparable to ours.
Weight quantization (GPTQ~\cite{frantar2023gptq}, AWQ~\cite{lin2024awq}, llama.cpp k-quants~\cite{kawrakow2023kquants}) and roofline analyses of decoding~\cite{williams2009roofline,pope2022scaling,yuan2024unveiled} explain why single-stream decoding is bandwidth-bound and why one weight read has to serve several tokens.

\paragraph{Exactness and determinism.}
Speculative decoding is lossless in exact arithmetic, but batched verification uses different kernels and reduction orders from single-token decoding, so greedy outputs can flip at near-ties.
vLLM states its guarantee only up to hardware numerical precision~\cite{vllm2026specdecode}.
Zhang et al.~\cite{zhang2026correctness} report 95\% exact match with standard decoding and trace the remainder to floating-point non-determinism.
An open llama.cpp issue~\cite{ankk2026issue25618} traces MTP greedy divergence on quantized targets to different single- and multi-token kernels, and restores bit-identity only by forcing the multi-token pass onto matching kernels.
Batch-invariant kernels~\cite{he2025nondeterminism,vllm2026batchinvariance} make results independent of batch size by fixing the reduction strategy, and LLM-42~\cite{gond2026llm42} enforces determinism with a verify-and-rollback loop over fixed-shape reductions.
Our kernel applies the batch-invariance principle to speculative verification on Metal.
It reproduces, row by row, the summation order of MLX's single-token \texttt{qmv} while still reading each weight once, so speculative output is byte-identical to plain greedy decoding, not only equal in distribution.
We found no prior work that makes a fast multi-token quantized verifier bit-identical to the single-token path; in our implementation the exactness costs 1--2\%.
